\section{Evidence-Calibrated Release}
\label{sec:evidence}

The theorem corpus and the release certificate answer different questions. A theorem is checked by the Lean kernel under its imports. A release certificate asserts that a named source subject has passed a specified collection of build and evidence gates. The evidence tuple is
\begin{equation}
E=\langle sha,tree,run,modules,cleanBuild,axioms,claims,
mutation,receipts,certificate,verdict\rangle.
\label{eq:evidence-tuple}
\end{equation}
The components must refer to one subject. A green job icon is not a substitute for inspecting their contents.

The formal inventory contains 145 declarations in the comprehensive audit and 27 declarations in the core-composition audit. The accepted axiom output is restricted to the standard dependencies
\begin{equation}
A_{allowed}=\{\mathsf{propext},\mathsf{Classical.choice},\mathsf{Quot.sound}\}.
\label{eq:allowed-axioms}
\end{equation}
No custom axiom is part of the audited proof boundary. These are commit-bound evidence facts; they must be regenerated after source changes.

Mutation testing supplies executable negative evidence. Let $b$ be a baseline fixture and $m_1,\ldots,m_k$ be property-breaking mutations. A genuine light gate requires
\begin{equation}
Gate(b,m_1,\ldots,m_k)=PASS
\iff Accept(b)\land\bigwedge_{i=1}^{k}\neg Accept(m_i).
\label{eq:mutation-gate}
\end{equation}
This criterion is \derived{} from the gate specification. It ensures that the report is generated by executions rather than fabricated as a static file. It does not establish completeness over all faults.

Cross-repository refinement receipts connect abstract verifier contracts to executable fixtures. For receipt (r), the required binding is
\begin{equation}
Binding(r)=\langle producerRepo,producerSha,consumerRepo,
consumerSha,task,digest,issued,expiry,authority\rangle.
\label{eq:cross-repo-binding}
\end{equation}
The ULM level formalizes narrower receipt components: one-rank authority succession and separate task, digest, time-window, and revocation predicates. Producer and consumer repository names and source SHAs in Equation~\eqref{eq:cross-repo-binding} are runtime-schema fields checked by the cross-repository pipeline, not fields proved by one ULM receipt theorem. A concrete receipt is therefore bounded runtime evidence. Equality between selected executable outputs and expected formal fixtures is not a universal proof of runtime equivalence.

Certificate generation is fail-closed. Missing required evidence produces a blocked status, generation or verification failure maps to job failure, and the independent verifier must agree with the certificate. Abstractly,
\begin{equation}
Verified(E)\iff SameSubject(E)\land CompleteEvidence(E)
\land CertificatePass(E)\land VerifierAgrees(E).
\label{eq:verified-release}
\end{equation}
This is a release-policy specification, not a Lean theorem. The machine-readable JSON artifacts, not this narrative, decide whether a particular subject satisfies it.

The release boundary also records what the formal development does not establish. Probability kernels, Bayesian calibration, generic analogy strength, a graph-similarity metric, human explanation quality, differential-privacy parameters, and substantive liability allocation are absent from the proved core. They may use the architecture's obligations and assurance types, but they require new semantics and evidence before receiving a \formalized{} label.

\subsection{Research integrity and reproducibility}

No private legal matter or customer dataset was used. Bibliographic entries were checked against DOI resolvers, publisher records, official journal sources, or EUR-Lex. The formal source remains primary for theorem claims. The recommended reproduction record is the exact commit, tree, workflow run, artifact names, certificate subject and status, missing-evidence array, and independent-verifier verdict.

The mutation gate and refinement receipts address different failure modes. Mutation asks whether a checker rejects deliberately invalid variants. Refinement asks whether an executable producer agrees with specified fixtures and binds its output to a source version. A system can pass one and fail the other. The certificate must therefore carry both results rather than replacing them with a single aggregate percentage.

Independent verification should recompute relations among artifacts, not merely trust their filenames. It checks subject identity, required fields, evidence status, and consistency between certificate and verifier verdict. When evidence is missing, the correct outcome is blocked rather than a partial pass disguised as success. This is fail-closed at the release layer in the same sense that `Outcome.Partial` and `solverIncomplete` remain explicit inside the model.

Evidence retention must also preserve negative results. Failed mutations, counterexamples to graph metrics, and the refuted privilege-to-privacy mapping are part of the scientific record. Deleting them would make later positive claims impossible to audit. Release documentation should state whether a limitation is unresolved, accepted as outside scope, or closed by a new theorem or fixture.

\subsection{Interpreting finite release evidence}

The evidence tuple should be read conjunctively and subject-specifically. A theorem audit establishes which declarations were checked under a particular source tree. A clean build establishes that the recorded environment compiled that tree. A mutation report establishes that enumerated property-breaking variants were rejected by the exercised checker. A refinement receipt establishes that selected executable fixtures produced the expected observations under named producer and consumer versions. The certificate and independent verifier then check consistency of those records. None of the components subsumes another, and none transfers automatically to a later commit.

Counts require the same discipline. A report that all designed mutations were rejected states coverage of the listed mutation catalogue, not coverage of all implementation defects. A report that three refinement fixtures passed states three witnessed executions, not three independent proofs or universal behavioral equivalence. The denominator must remain visible because ``all passed'' without the fixture inventory is not reproducible. Similarly, the 145- and 27-declaration audits are inventory facts for the reported subject; they are not permanent measures of project size or theorem strength.

Evidence can fail in several non-equivalent ways. A missing artifact is incompleteness. A subject mismatch is invalid composition. A failed mutation shows that the exercised checker does not enforce an intended property. A refinement mismatch shows disagreement between selected executable and reference observations. A verifier disagreement shows that the certificate interpretation is unstable. Fail-closed reporting should preserve these diagnoses, since each has a different repair. Regenerating a filename does not repair a failed fixture, and rerunning a build does not repair a source-identity mismatch.

Archival practice is part of the scholarly claim. An evidence citation should resolve to immutable or content-addressed artifacts containing the subject SHA, run identity, report schema version, fixture identifiers, certificate status, missing-evidence list, and verifier verdict. If artifacts remain only in an expiring CI service, the paper can describe the method but cannot independently support exact historical counts from the checkout alone. A repository release, durable archive, or deposited bundle should therefore accompany any numerical release statement.

The independent verifier should be operationally separate enough to catch generator mistakes, but independence is not absolute. It may share parsers, schemas, libraries, or assumptions with the generator. The assurance case should record that trusted computing base and avoid treating a second JSON verdict as epistemically independent merely because it has a different filename. Mutation of certificate fields, cross-subject fixtures, and intentionally incomplete evidence sets are useful negative cases for the verifier itself.

Finally, release is revocable in the ordinary scientific sense. Discovery of a missing premise, incorrect fixture, compromised toolchain, or misidentified subject does not make the earlier artifact disappear; it changes the current assessment of its force. A superseding record should retain the prior evidence, state the reason for withdrawal, and bind the correction to a new subject. This preserves the audit trail and prevents later readers from confusing historical execution success with a standing warranty of legal or software correctness.
